% SLICEBONFERRONI1 2026-10-16 08:40 Asia/Seoul
\documentclass[11pt]{article}
\usepackage[T1]{fontenc}
\usepackage{lmodern,amsmath,amssymb,amsthm,booktabs,hyperref}
\usepackage[margin=0.9in]{geometry}
\newtheorem{theorem}{Theorem}
\newtheorem{lemma}[theorem]{Lemma}
\newtheorem{corollary}[theorem]{Corollary}
\newcommand{\EXACT}{\textbf{PROVED/EXACT}}
\newcommand{\KNOWN}{\textbf{KNOWN THEOREM}}
\newcommand{\COMP}{\textbf{COMPUTATIONAL}}
\newcommand{\OPEN}{\textbf{OPEN}}
\newcommand{\RETIRED}{\textbf{FALSIFIED/RETIRED}}
\title{SLICEBONFERRONI1:\\Same-Layer Sunflower Obstruction and a Moment--Bonferroni Bridge}
\author{P versus NP project}
\date{16 October 2026}
\begin{document}
\maketitle

\begin{center}\fbox{\parbox{0.94\linewidth}{
\textbf{Status.} The same-layer obstruction to direct application of the Cavalar--Fabris--Mukhopadhyay--Srinivasan--Yehudayoff lower-bound criterion, the Bonferroni moment-fooling lemma, and the exact negative overlap functional are \EXACT\ at their stated models, pending independent review. The finite arithmetic audit is \COMP. The circuit-to-low-overlap-DNF approximation theorem and the asymptotic UNSAT moment-matching theorem remain \OPEN. Therefore $P\ne NP$ remains \OPEN.
}}\end{center}

\section{IDEAS PROPOSED BEFORE LITERATURE}
\paragraph{CANDIDATE A: direct same-layer robust-sunflower criterion.}
Use the exact-$3q$ UNSAT distribution as accepting inputs and the uniform exact-$3q$ slice as rejecting test distribution in a modern distributional monotone lower-bound theorem. Falsifier: fixed Hamming weight forces accepting spread too small relative to the rejecting robust-sunflower threshold.

\paragraph{CANDIDATE B: moment--Bonferroni fooling.}
After a small monotone circuit is approximated by a bounded-width DNF, match sufficiently many clause moments between the UNSAT positive distribution and the uniform slice. Use truncated inclusion--exclusion rather than generic positive spread to prove the approximator cannot distinguish them. Missing theorem: control of the Bonferroni remainder for every allowed approximator.

\paragraph{CANDIDATE C: signed-DAG union spectrum.}
Use the previous signed-hypergraph canonicalization only to measure the union-size spectrum of approximator terms, not as a raw state-count lower bound. Prediction: closure/sunflower pruning forces enough union growth that high simultaneous term multiplicity is rare on the uniform slice.

\paragraph{CANDIDATE D: HD/orbit moment search.}
Use the orbit-moment polytope and proof-safe HD sketches from SLICEOUTPUTMOMENT1 to search for exact moment-matching distributions through the order demanded by Candidate B. Every projected certificate must be lifted or accompanied by an injectivity certificate.

\section{LITERATURE CHECK AND REVISIONS}
Cavalar et al.\ give a general lower-bound criterion (their Theorem 3.3): an accepting distribution must be $t$-wise $Q$-spread, a rejecting distribution must have robust-sunflower parameter $r_w$, and the criterion requires $8r_w\le Q$ before yielding a monotone circuit lower bound. This is a direct model match for Candidate A. The exact same-layer calculation below shows that this condition cannot hold for our uniform exact-layer rejecting distribution once the accepting distribution is also supported on that exact layer.

Cavalar--Kumar--Rossman establish gate-by-gate monotone approximation using closure and robust sunflowers. The 2026 derandomized sunflower paper states that robust-sunflower lemmas extend to $O(w^2)$-wise independent distributions, and develops further method-of-moments relaxations. Bazzi proves that $O(\log^2 m)$-wise independence fools polynomial-size DNF/CNF formulas under the cube distribution. These results support Candidate B, but none directly gives a fixed-slice DNF-fooling theorem for our positive distribution.

Bonferroni inequalities are precisely truncated inclusion--exclusion. This suggests a direct self-contained bridge from exact low-order moment matching to fooling a particular DNF, provided simultaneous satisfaction of too many terms has small factorial moment. Candidate B is therefore developed below.

The recent spread-matching and tailored-sunflower monotone lower bounds show that problem-specific overlap structure can outperform generic spread arguments. Candidate C is retained as the next combinatorial target. No theorem from that literature is silently transferred to our SAT slice.

\section{CONSTRUCTION OR PROOF ATTEMPT}

\subsection{Exact same-layer obstruction to Theorem 3.3}
Let the clause-coordinate universe have size
\[
N=8\binom q3,\qquad t=3q,\qquad p=t/N.
\]
Let $\mu_1$ be \emph{any} distribution supported on the exact layer $\binom{[N]}t$, and let $\mu_0$ be uniform on that exact layer. Let $Q$ denote the spread parameter in the convention
\[
\Pr_{\mu_1}[A\subseteq X]\le Q^{-|A|}.
\]

\begin{lemma}[Accepting spread ceiling]
Every such $\mu_1$ satisfies
\[
Q\le N/t.
\]
\end{lemma}
\begin{proof}
The singleton marginals sum to $t$. Hence some coordinate has marginal at least $t/N$. The $Q$-spread inequality on that singleton gives $Q^{-1}\ge t/N$.
\end{proof}

For a distribution $D$, let $r(D,\ell,\varepsilon)$ be the minimum $r$ such that every $\ell$-uniform family of at least $r^\ell$ terms contains a $(D,\varepsilon)$-robust sunflower, in the convention of Cavalar et al.

\begin{lemma}[Uniform-slice singleton sunflower floor]
Assume $N/t>4$ and $0<\varepsilon<1/2$. Then
\[
r(\mu_0,1,\varepsilon)>
\left\lfloor\frac{N-t+1}{2t}\right\rfloor.
\]
\end{lemma}
\begin{proof}
Let
\[
m=\left\lfloor\frac{N-t+1}{2t}\right\rfloor
\]
and take a family of $m$ distinct singleton terms. Its core is empty. The probability that none of these $m$ coordinates is selected by a uniform $t$-subset is
\[
\frac{\binom{N-m}{t}}{\binom Nt}
=\prod_{i=0}^{t-1}\left(1-\frac{m}{N-i}\right)
\ge
\left(1-\frac1{2t}\right)^t
\ge\frac12,
\]
where the last inequality is Bernoulli's inequality. Thus the probability that at least one petal is satisfied is at most $1/2<1-\varepsilon$. This family is not a $(\mu_0,\varepsilon)$-robust sunflower.
\end{proof}

\begin{theorem}[Same-layer criterion obstruction]
For all sufficiently large $q$, no choice of exact-layer accepting distribution $\mu_1$ can satisfy the hypothesis
\[
8r_w\le Q
\]
of Theorem 3.3 of Cavalar et al.\ when the rejecting test distribution is the uniform exact-$3q$ layer and $r_w$ includes the $\ell=1$ robust-sunflower parameter at any error $\varepsilon<1/2$.
\end{theorem}
\begin{proof}
The previous lemmas give $Q\le N/t$ and
\[
r_w\ge r(\mu_0,1,\varepsilon)>
\left\lfloor\frac{N-t+1}{2t}\right\rfloor.
\]
For $N/t\to\infty$, the latter is $(1/2-o(1))N/t$, hence
\[
8r_w>(4-o(1))N/t>Q.
\]
Therefore the required inequality is impossible.
\end{proof}

This retires Candidate A at the exact declared interface. It does \emph{not} rule out all robust-sunflower or approximation arguments; it shows that the modern generic spread-vs-sunflower criterion cannot be applied directly with both test distributions on the same fixed layer.

\subsection{Moment--Bonferroni bridge for bounded-width DNFs}
Let
\[
D(x)=\bigvee_{i=1}^{L} t_{T_i}(x),\qquad
t_T(x)=\prod_{a\in T}x_a,
\]
where $|T_i|\le w$. Define
\[
R(x)=\sum_{i=1}^{L}t_{T_i}(x),
\]
the number of DNF terms satisfied by $x$.

\begin{lemma}[Pointwise Bonferroni remainder]
For every integer $R\ge1$ and $s\ge1$,
\[
1-\sum_{j=1}^{s}(-1)^{j+1}\binom Rj
=(-1)^s\binom{R-1}{s}.
\]
Moreover
\[
\binom{R-1}{s}\le \binom R{s+1}
\]
whenever the left side is nonzero.
\end{lemma}
\begin{proof}
The first identity is the standard alternating binomial identity. If $R\le s$, both remainders vanish. If $R\ge s+1$,
\[
\binom{R-1}s=\frac{s+1}{R}\binom R{s+1}\le\binom R{s+1}.
\]
\end{proof}

\begin{theorem}[Exact moment--Bonferroni DNF fooling]
Let $\mu,\nu$ be distributions on $\{0,1\}^N$ that have identical coordinate-conjunction moments
\[
\mathbb E_\mu[t_A]=\mathbb E_\nu[t_A]
\]
for every $A$ of size at most $(s+1)w$. Then
\[
\left|\Pr_\mu[D=1]-\Pr_\nu[D=1]\right|
\le
2\,\mathbb E_\nu\binom R{s+1}.
\]
The same bound holds with $\mu$ in place of $\nu$ on the right.
\end{theorem}
\begin{proof}
For every $j\le s+1$,
\[
\binom Rj
=
\sum_{\substack{J\subseteq[L]\\|J|=j}}
t_{\cup_{i\in J}T_i}.
\]
Each union contains at most $jw\le(s+1)w$ coordinates. Hence moment matching implies equality of the expectations of $\binom Rj$ under $\mu$ and $\nu$ for all $j\le s+1$.

Apply the pointwise Bonferroni identity to $1_{R\ge1}$. The expectations of the truncated sums through order $s$ cancel between $\mu$ and $\nu$. The two absolute remainders are bounded by the corresponding expectations of $\binom R{s+1}$, and those expectations are equal by the preceding paragraph.
\end{proof}

This theorem supplies the previously missing exact connection between orbit-moment matching and a concrete approximator class. It does not require generic positive spread.

\subsection{Exact negative overlap functional}
For the uniform exact-$t$ slice $\nu$, a fixed coordinate set $A$ of size $a$ is contained with probability
\[
h_{N,t}(a)=\frac{(t)_a}{(N)_a}.
\]
Therefore
\[
\mathbb E_\nu\binom R{s+1}
=
\sum_{\substack{J\subseteq[L]\\|J|=s+1}}
h_{N,t}\!\left(\left|\bigcup_{i\in J}T_i\right|\right).
\]
Define the exact overlap functional
\[
\Xi_{s+1}(\mathcal T;N,t)
:=
\sum_{\substack{J\subseteq[L]\\|J|=s+1}}
\frac{(t)_{|\cup_{i\in J}T_i|}}{(N)_{|\cup_{i\in J}T_i|}}.
\]
Then exact moment matching through order $(s+1)w$ gives
\[
\boxed{
\left|\Pr_\mu[D=1]-\Pr_\nu[D=1]\right|
\le 2\Xi_{s+1}(\mathcal T;N,t).
}
\]

The previous signed-hypergraph DAG work can now be used in a sharply limited way: canonicalize the term family and record the distribution of union sizes of $(s+1)$-tuples. Raw DAG size is irrelevant; only the induced value of $\Xi_{s+1}$ matters.

A simple sufficient condition is
\[
\left|\bigcup_{i\in J}T_i\right|\ge u
\quad\text{for every }|J|=s+1,
\]
which implies
\[
\Xi_{s+1}\le \binom L{s+1}(t/N)^u.
\]
Thus the next combinatorial objective is a \emph{union-growth} theorem for the term families produced by closure/sunflower pruning, or a direct average bound on $\Xi$.

\subsection{New proof spine}
For a fixed polynomial circuit exponent $K$, it would suffice to prove all three statements below.

\begin{enumerate}
\item \textbf{Approximation.} Every size-$N^K$ monotone circuit is approximated on both test distributions by a monotone DNF $\mathcal T$ of width
\[
w_K=O_K(\log q)
\]
with total approximation error $o(1)$.
\item \textbf{Moment design.} There is an exact-$3q$ UNSAT distribution matching the uniform slice on all coordinate moments through order
\[
r_K=(s_K+1)w_K
\]
for some $s_K$ with $r_K=O_K(\log q)$.
\item \textbf{Overlap control.} Every resulting approximator satisfies
\[
\Xi_{s_K+1}(\mathcal T;N,3q)=o(1).
\]
\end{enumerate}

Then the approximator cannot distinguish the positive and negative distributions by the moment--Bonferroni theorem, while a correct approximator to the exact slice must distinguish them with advantage $1-o(1)$, a contradiction. Valiant's slice theorem would then transfer the monotone lower bound to unrestricted circuit complexity.

The point of this formulation is that the moment order is only logarithmic if $s_K$ is constant and $w_K=O_K(\log q)$. That is precisely the regime where SLICEOUTPUTMOMENT1's signed-orbit space has dimension
\[
q^{O_K(\log\log q)}
\]
and where the high-girth planting is locally invisible.

\section{TESTS AND OUTCOME}
A bounded arithmetic audit verified the Bonferroni identity and remainder inequality for $1\le s\le8$ and $0\le R\le30$ (248 cases, zero discrepancies). For the exact density-3 clause universe, the singleton robust-sunflower obstruction was evaluated at $q=8,16,32,64,128$. The constructed non-robust singleton families had no-hit probabilities $0.6414,0.6077,0.6064,0.6064,0.6064$, respectively, and in every case $8m>N/t$. These finite checks are \COMP; the theorems above are symbolic.

\paragraph{Outcome.}
Candidate A is \RETIRED\ at the declared direct same-layer criterion. Candidate B is \EXACT\ as a DNF-fooling lemma but does not yet cover arbitrary small monotone circuits. Candidate C is the highest-value new combinatorial target: prove that closure/sunflower-pruned approximator term families have $\Xi_{s+1}=o(1)$. Candidate D remains a search/certificate layer for the moment-design obligation.

\paragraph{Next target.}
Derive an overlap/union-growth theorem from the absence of a robust sunflower or from a tailored spread-matching condition. In parallel, use the signed-orbit DAG to enumerate the worst union spectra of small closed term families, looking for the extremal configuration that maximizes $\Xi_{s+1}$. This is now a much narrower question than the original monotone lower bound.

Root \texttt{main.tex} remains truncated and is not modified.

\begin{thebibliography}{6}
\bibitem{CFMSY25} B. P. Cavalar, T. Fabris, P. Mukhopadhyay, S. Srinivasan, A. Yehudayoff, \emph{Negations Are Powerful Even in Small Depth}, ECCC TR25-219 / STOC 2026.
\bibitem{CKR22} B. P. Cavalar, M. Kumar, B. Rossman, \emph{Monotone Circuit Lower Bounds from Robust Sunflowers}, Algorithmica 84, 2022.
\bibitem{CFMSY26} B. P. Cavalar, T. Fabris, P. Mukhopadhyay, S. Srinivasan, A. Yehudayoff, \emph{Derandomized Sunflowers and Radical Lower Bounds}, ECCC TR26-220, 2026.
\bibitem{Bazzi09} L. Bazzi, \emph{Polylogarithmic Independence Can Fool DNF Formulas}, SIAM J. Comput. 38(6), 2009.
\bibitem{Rao26} A. Rao, \emph{Monotone Circuit Lower Bounds from Spread Matchings}, ECCC TR26-129, 2026.
\bibitem{Valiant86} L. G. Valiant, \emph{Negation Is Powerless for Boolean Slice Functions}, SIAM J. Comput. 15(2), 1986.
\end{thebibliography}
\end{document}
% FINAL_VERIFICATION_SENTINEL::MAIN::SLICEBONFERRONI1::2026-10-16_0840::P_VS_NP_OPEN
